\section*{Abstract}
\addcontentsline{toc}{section}{Abstract}

This report presents a short-term electricity demand forecasting study for
Tetouan, Morocco. The operating problem is to anticipate city load one hour
ahead, and to raise an alert when the forecast crosses a safety capacity
threshold, so that reserve margin can be armed before the evening peak. The
source is the UCI \emph{Power Consumption of Tetouan City} series: 10-minute
measurements from three distribution zones, together with temperature,
humidity, wind speed and diffuse solar radiation. The series is aggregated to
hourly demand. Calendar variables, lagged demand, rolling statistics and
weather features feed a seasonal-naive baseline, a ridge regression, a random
forest, and two gradient-boosting models (XGBoost and LightGBM). Models are
compared on a future test block with MAE, RMSE and WAPE, and peak risk is
scored with peak-magnitude error, peak-timing error and breach anticipation
rate. Square-bracket entries in the result tables are completed from
\path{source_code/outputs/metrics.json} after the pipeline is run.

\noindent\textbf{Keywords:} electricity demand; short-term load forecasting;
time series; gradient boosting; peak capacity; Tetouan.
